% 04_VRL_Kvantifikatory.tex
% Příprava: 1. ročník, hodina 04, Výroková logika; běžná hodina 45 minut.
% Návaznost: 02_VRL_CoJeVyrok.tex (otevřená formule),
% 03_VRL_PravdivostSlozenychVyroku.tex (negace, implikace, pravdivost).
% Plány: 1_rocnik/tematicke_plany/SPS.pdf a GYM.pdf, oddíl 1.1.
% Obsahový podklad: ucebnice/01_Zakladni_poznatky/Vyrokova_logika/
% Výroková logika.tex, oddíly „Kvantifikované výroky a jejich negace“
% a „Logická výstavba matematiky“; vlastní formulace výkladu.
% Procvičení a): vzor z úlohy „Kvantifikátory ve větách“, druhá položka;
% b): „Negace výroků – jednoduché“, položka „Každé prvočíslo je liché“;
% c): „Negace výroků s kvantifikátory“, poslední položka.
% Ostatní konkrétní zadání a řešení: vlastní tvorba pro tuto přípravu.
% Relevantní místní podklad Realisticky.cz nebyl nalezen; není použit.
% Kompilace z Work/MAT: latexmk -cd -pdf -interaction=nonstopmode -halt-on-error 1_rocnik/pripravy/04_VRL_Kvantifikatory.tex
\documentclass[11pt,a4paper]{article}
% Spolecny styl pro ucitelske pripravy MAT (Work/MAT)
% Verze 1.1, 2026-10-23
\NeedsTeXFormat{LaTeX2e}
\ProvidesPackage{priprava}[2026/10/23 v1.1 Sablona ucitelskych priprav MAT]

% Zakladni typografie a cestina
\RequirePackage[T1]{fontenc}
\RequirePackage[utf8]{inputenc}
\RequirePackage[czech]{babel}
\RequirePackage{lmodern}
\RequirePackage{microtype}
\RequirePackage{amsmath,amssymb}
\RequirePackage{xcolor}
\RequirePackage{enumitem}
\RequirePackage[most]{tcolorbox}
\RequirePackage{tabularx,booktabs,array}
\RequirePackage{fancyhdr}
\RequirePackage{needspace}
\RequirePackage{geometry}
\RequirePackage[hidelinks]{hyperref}

\geometry{a4paper,top=17mm,bottom=18mm,left=23mm,right=23mm}
\definecolor{notegray}{gray}{0.39}
\setlength{\parindent}{0pt}
\setlength{\parskip}{2.5pt plus 1pt minus 0.5pt}
\setlist[itemize]{topsep=2pt,itemsep=1pt,parsep=0pt,leftmargin=1.45em}
\setlist[enumerate]{topsep=2pt,itemsep=2pt,parsep=0pt,leftmargin=1.65em}
\setcounter{secnumdepth}{0}
\renewcommand{\arraystretch}{1.12}

% Zahlavi: vlevo rocnik | globalni tema, vpravo cislo: nazev hodiny
% Zapati zustava zcela prazdne.
\pagestyle{fancy}
\fancyhf{}
\renewcommand{\headrulewidth}{0.35pt}
\renewcommand{\footrulewidth}{0pt}
\setlength{\headheight}{14pt}
\newcommand{\nastavPripravu}[4]{%
  \fancyhead[L]{\small #1\enspace\textbar\enspace #2}%
  \fancyhead[R]{\small #3: #4}%
}

% Minimalisticke cernobile boxy. Kazda definice zvlast.
\tcbset{commonlesson/.style={enhanced,breakable=false,colback=white,colframe=black,
  colbacktitle=white,coltitle=black,boxrule=0.45pt,
  arc=0.8mm,outer arc=0.8mm,left=3mm,right=3mm,top=1.6mm,bottom=1.6mm,
  toptitle=1mm,bottomtitle=0.7mm,fonttitle=\bfseries,
  attach boxed title to top left={xshift=2.5mm,yshift=-1.15mm},
  boxed title style={colback=white,colframe=white,boxrule=0pt,left=1mm,right=1mm}}}
\newtcolorbox{zadaniBox}[1][]{commonlesson,title={Zadání},#1}
\newtcolorbox{prikladBox}[1][]{commonlesson,title={Příklad},#1}
\newtcolorbox{definiceBox}[1][]{commonlesson,title={Definice},#1}

% Spolecny box pro zadani a jeho reseni, bez mezery mezi castmi.
% Pouziti:
%   \begin{prikladReseniBox}
%     Zadani ulohy.
%     \tcblower
%     \begin{align*} ... \end{align*}
%   \end{prikladReseniBox}
% Spodni cast ma jemny sedy podklad a delici linku; lze v ni sazet
% vice radku matematiky. Obe casti maji jen jeden spolecny ramecek.
\tcbset{solutionlesson/.style={
  commonlesson,
  enhanced jigsaw,
  breakable,
  colbacklower=black!3,
  segmentation style={solid,draw=black!25,line width=0.35pt},
  fontlower=\small\color{notegray},
  before lower={\noindent\textbf{Řešení:}\enspace},
  leftlower=3mm,rightlower=3mm
}}
\newtcolorbox{prikladReseniBox}[1][]{solutionlesson,title={Příklad},#1}
\newtcolorbox{zadaniReseniBox}[1][]{solutionlesson,title={Zadání},#1}

% Pomucky pro autora pripravy: sedive poznamky; \reseni pro starsi pripravy
\newcommand{\poznamka}[1]{\par\begingroup\small\color{notegray}\itshape
  \textnormal{\textbf{Poznámka pro učitele.}}\enspace #1\par\endgroup}
\newcommand{\reseni}[1]{\par\begingroup\small\color{notegray}
  \textbf{Očekávané řešení:}\enspace #1\par\endgroup}
\newcommand{\blok}[1]{\needspace{5\baselineskip}\subsection*{#1}}
\newcommand{\vyrok}[1]{\uv{#1}}

\nastavPripravu{1. ročník}{Výroková logika}{04}{Kvantifikátory}

\begin{document}

\blok{1. Od otevřené formule k výroku}
Vrať se k otázce z druhé hodiny: je zápis $x^2=1$ s neurčeným $x$
výrokem? Pro $x=1$ rovnost platí, pro $x=0$ neplatí. Minule jsme
určovali pravdivost složených výroků; nyní upřesníme, \emph{pro které
hodnoty proměnné} má tvrzení platit.

\begin{definiceBox}
\textbf{Výroková forma (otevřená formule)} $P(x)$ je sdělení
s proměnnou $x$ z určené množiny $M$, které se po dosazení konkrétního
prvku za $x$ stane výrokem.
\end{definiceBox}

Výrok získáme také tím, že řekneme, zda vlastnost požadujeme pro všechny
prvky dané množiny, nebo alespoň pro jeden. K tomu slouží kvantifikátory.
Zápis $x\in M$ čteme \uv{$x$ je prvkem množiny $M$}.

\blok{2. Pro každé a existuje}
\begin{definiceBox}
\textbf{Obecný (univerzální) kvantifikátor} $\forall$ vyjadřuje
\uv{pro každé}. Výrok $\forall x\in M:\ P(x)$ je pravdivý právě tehdy,
když $P(x)$ platí pro každý prvek množiny $M$.
\end{definiceBox}

\begin{definiceBox}
\textbf{Existenční kvantifikátor} $\exists$ vyjadřuje
\uv{existuje alespoň jeden}. Výrok $\exists x\in M:\ P(x)$ je pravdivý
právě tehdy, když v množině $M$ existuje alespoň jeden prvek, pro který
$P(x)$ platí.
\end{definiceBox}

\begin{prikladReseniBox}
Nechť $M=\{-2,-1,0,1,2\}$. Přečti a rozhodni o pravdivosti výroků:
\[
 A:\ \forall x\in M:\ x^2=1,
 \qquad
 B:\ \exists x\in M:\ x^2=1.
\]
\tcblower
$A$: \uv{Pro každé číslo $x$ z $M$ platí $x^2=1$.} Výrok je nepravdivý,
například pro $x=0$ dostáváme $0^2=0\ne1$.

$B$: \uv{Existuje alespoň jedno číslo $x$ z $M$, pro které platí $x^2=1$.}
Výrok je pravdivý: vyhovují $x=-1$ i $x=1$.
\end{prikladReseniBox}

K vyvrácení obecného výroku stačí jediný nevyhovující prvek.
K potvrzení existenčního výroku stačí jediný vyhovující prvek.
Několik vyhovujících příkladů však obecný výrok o nekonečné množině
nedokazuje.

\poznamka{Zdůrazni, že \uv{existuje} neznamená \uv{existuje právě jeden}.
Množina, z níž vybíráme, patří k zadání: v množině $\{0,2\}$ by výrok
$\exists x:\ x^2=1$ neplatil. Nezaváděj nyní více kvantifikátorů za sebou.}

\newpage
\blok{3. Matematické věty}
\begin{definiceBox}
\textbf{Matematická věta} je matematické tvrzení, jehož pravdivost
je dokázána na základě přijatých výchozích tvrzení (axiomů), definic
a již dokázaných vět.
\end{definiceBox}

Matematické věty často vyjadřují obecnou platnost. Slovo \uv{každé}
bývá skryto i ve větách tvaru \uv{Je-li splněn předpoklad, pak platí závěr.}
Předpoklady i množina, o níž mluvíme, musí být zřejmé.

\begin{prikladReseniBox}
Zapiš pomocí kvantifikátoru větu \uv{Druhá mocnina každého reálného
čísla je nezáporná} a zdůvodni ji pro libovolné reálné číslo.
\tcblower
Zápis je $\forall x\in\mathbb{R}:\ x^2\geq0$.
Je-li $x=0$, pak $x^2=0$. Je-li $x\ne0$, jsou oba činitelé
v součinu $x\cdot x$ stejného znaménka, a součin je proto kladný.
Pro všechna reálná čísla tedy platí $x^2\geq0$.
\end{prikladReseniBox}

\poznamka{Značka $\mathbb{R}$ označuje reálná čísla. Porovnej obecné
zdůvodnění s pouhým dosazením několika čísel. Systematické metody důkazů
nejsou náplní této hodiny.}

\blok{4. Negace kvantifikovaných výroků}
Z minulé hodiny víme, že negace má opačnou pravdivostní hodnotu.
Věta \uv{Všichni žáci dané třídy odevzdali úkol} neplatí, jakmile
\emph{alespoň jeden} žák úkol neodevzdal. Její negace tedy neříká,
že úkol neodevzdal nikdo.

Naopak výrok \uv{Alespoň jeden žák dané třídy odevzdal úkol}
je nepravdivý, právě když úkol neodevzdal \emph{žádný} žák této třídy.
Při negaci proto zaměníme kvantifikátor a současně negujeme vlastnost:
\begin{align*}
 \neg\bigl(\forall x\in M:\ P(x)\bigr)
   &\Leftrightarrow \exists x\in M:\ \neg P(x),\\
 \neg\bigl(\exists x\in M:\ P(x)\bigr)
   &\Leftrightarrow \forall x\in M:\ \neg P(x).
\end{align*}
Množina $M$ zůstává stejná. Negaci tvoříme i u nepravdivého výroku.

\begin{prikladReseniBox}
Utvoř negace následujících výroků, přečti je slovy a urči pravdivost
původních výroků i jejich negací:
\begin{enumerate}[label=\alph*)]
  \item $\forall x\in\mathbb{R}:\ x^2\geq0$.
  \item $\exists x\in\mathbb{R}:\ x^2=-1$.
\end{enumerate}
\tcblower
\begin{enumerate}[label=\alph*)]
  \item $\exists x\in\mathbb{R}:\ x^2<0$.
  \uv{Existuje reálné číslo, jehož druhá mocnina je záporná.}
  Původní výrok je pravdivý, negace nepravdivá, jak plyne z uvedené věty.
  \item $\forall x\in\mathbb{R}:\ x^2\ne-1$.
  \uv{Žádné reálné číslo nemá druhou mocninu rovnou $-1$.}
  Protože $x^2\geq0$, je původní výrok nepravdivý a negace pravdivá.
\end{enumerate}
\end{prikladReseniBox}

\poznamka{Při negování nerovnosti hlídej rovnost na hranici:
negací $u\geq0$ je $u<0$, negací $u>0$ je $u\leq0$.
Nestačí jen zaměnit $\forall$ a $\exists$ nebo jen přidat zápor.}

\newpage
\blok{5. Samostatné procvičení}
\begin{zadaniReseniBox}
U každé položky zapiš původní výrok i jeho negaci pomocí kvantifikátoru,
negaci přečti slovy a urči pravdivost obou výroků. Odpověď zdůvodni.
\begin{enumerate}[label=\alph*)]
  \item Doplň kvantifikátor tak, aby vznikl pravdivý výrok o všech
  reálných číslech: $\ldots\ x\in\mathbb{R}:\ x^2+1>0$.
  \item Každé prvočíslo je liché.
  \item Existuje přirozené číslo, které je menší než nula.
\end{enumerate}
\tcblower
\begin{enumerate}[label=\alph*)]
  \item Původní výrok: $\forall x\in\mathbb{R}:\ x^2+1>0$.
  Je pravdivý, neboť $x^2\geq0$, tedy $x^2+1\geq1>0$.
  Negace: $\exists x\in\mathbb{R}:\ x^2+1\leq0$.
  \uv{Existuje reálné číslo, pro které je $x^2+1$ nekladné.}
  Negace je nepravdivá.
  \item Označíme-li $M$ množinu všech prvočísel, je původní výrok
  $\forall p\in M:\ p\text{ je liché}$.
  Negace: $\exists p\in M:\ p\text{ není liché}$, tedy
  \uv{Existuje prvočíslo, které je sudé.}
  Číslo $2$ je prvočíslo a je sudé; původní výrok je nepravdivý,
  negace pravdivá.
  \item Původní výrok: $\exists n\in\mathbb{N}:\ n<0$.
  Negace: $\forall n\in\mathbb{N}:\ n\geq0$.
  \uv{Každé přirozené číslo je nezáporné.}
  Původní výrok je nepravdivý, negace pravdivá.
\end{enumerate}
\end{zadaniReseniBox}

\poznamka{Prvočíslo je přirozené číslo větší než $1$ s právě dvěma
kladnými děliteli: $1$ a sebou samým. Značka $\mathbb{N}$ označuje
přirozená čísla; výsledek c) nezávisí na tom, zda mezi ně zahrnujeme nulu.
Při pomalejším tempu řeš společně a), samostatně b) a c);
následující rozšíření vynech.}

\blok{Volitelné rozšíření: kvantifikátor a implikace}
\begin{prikladReseniBox}
Neguj výrok $\forall n\in\mathbb{Z}:\ (n>0\Rightarrow n>1)$.
Urči pravdivost původního výroku i negace.
\tcblower
Z minulé hodiny: implikace $p\Rightarrow q$ je nepravdivá právě tehdy,
když $p$ platí a $q$ neplatí. Proto
\begin{align*}
 \neg\bigl(\forall n\in\mathbb{Z}:\ (n>0\Rightarrow n>1)\bigr)
 &\Leftrightarrow \exists n\in\mathbb{Z}:\ (n>0\land n\leq1).
\end{align*}
Slovy: \uv{Existuje celé číslo, které je kladné a zároveň není větší
než jedna.} Vyhovuje $n=1$. Původní výrok je nepravdivý a negace pravdivá.
\end{prikladReseniBox}

\poznamka{Značka $\mathbb{Z}$ označuje celá čísla. Negací implikace
není její obrácení. Například $n=0$ původní obecný výrok nevyvrací:
předpoklad $n>0$ je pro něj nepravdivý.}

\textbf{Shrnutí na tabuli:} $\forall$: pro každé.
$\exists$: alespoň jeden. Při negaci zaměníme kvantifikátor,
negujeme vlastnost a zachováme množinu.

\end{document}
